\chapter{Der Tod ein König der Schrecken}
\quellseite{10}

Vor zwei Jahren starb in einer Straße der inneren Stadt (Berlin) ein Schuhmacher, der sein Lebenlang ein wüstes und ungöttliches Leben geführt hatte. Seine Ehe war nicht eingesegnet und die Kinder waren nicht getauft. Aus der Kirche war er ausgetreten. Seine Kinder wußten ebensowenig wie er selbst und seine Frau von Religion. Nun aber kam die Todesstunde. Die gibt den Menschen fast immer auf, etwas zu denken, und ihm auch. Als es ans Sterben ging und er in den letzten Tagen halb bewußtlos dalag, verlangte er von seiner Frau, sie solle mit ihm beten, aber sie konnte es nicht. Dann verlangte er von seinen Kindern, sie möchten mit ihm das „Vater-Unser“ beten, aber keins von seinen fünf Kindern konnte das. Da sollten sie ein Gesangbuch oder ein Neues Testament suchen, aber trotz allen Suchens fanden sie keins. Und schließlich verlangte er nach einem Prediger, aber niemand von seinen Angehörigen kannte einen solchen. So ist er in seiner Verzweiflung dahingefahren, als ein erschütterndes Beispiel dafür, daß der Tod den Menschen viel zu denken gibt.

(Blätter aus der Stadtmission.)
